\section{Related Work}
\label{sec:related}

\paragraph{Language models as models of the brain and of reading.}
Hidden states of language models predict fMRI, MEG and intracranial responses during naturalistic comprehension \citep{schrimpf2021neural,caucheteux2022brains,caucheteux2023evidence,toneva2019interpreting}, and this alignment generally improves with model quality and size, in line with the Platonic Representation Hypothesis \citep{huh2024platonic}; it also emerges in models trained on child-scale corpora \citep{hosseini2024artificial} and saturates early in training \citep{alkhamissi-etal-2025-language}. Surprisal, in turn, is the dominant predictor of reading times \citep{hale2001probabilistic,levy2008expectation,smithlevy2013effect,wilcox2023testing,shain2024large}, but here larger and better models are often \emph{worse} predictors \citep{ohschuler2023why,ohschuler2023twobillion,kuribayashi2021lower,wilcox2024bigger}. Concurrently, \citet{devarda2026behavioral} show within GPT-2 that a low-dimensional effort predictor suffices for behavior, whereas brain responses need rich embeddings, and that smaller models fit behavior and larger ones the brain. We hold the predictors fixed and ask the complementary question: which \emph{model class} is closer to each target.

\paragraph{Memory and developmental constraints.}
Architectural constraints restrict what context a model can use: capacity-limited or single-head attention \citep{timkey2023limited}, recency biases such as ALiBi \citep{press2021train,devarda2024locally,clark2025linear,thamma2025humanlike}, forget gates \citep{lin2025forgetting}, stick-breaking attention \citep{tan2025scaling}, static local windows \citep{bu2026staticlinearattentionconstraints}, and resource-rational or memory-bounded processing \citep{hahn2022resourcerational,xu2026memory,xufutrell2026strategic,keshev2025working}. They are motivated by retrieval-based and lossy-context theories of sentence processing \citep{lewisvasishth2005activation,futrell2020lossy,gibson2000dlt}. Lossy-context surprisal has recently been tested at the level of individual readers: readers with lower working-memory capacity show locality effects at shorter dependencies, as if they retained a smaller share of the preceding context, and their reading times follow lossy-context surprisal more closely \citep{li2027modeling}. The BabyLM Challenge keeps the architecture and restricts the training data \citep{warstadt2023cfp,warstadt2023findings,hu2024secondbabylm,huebner2021babyberta}. These approaches are usually evaluated on one kind of human data and against a few baselines; we place them in one pipeline, cross them as two factors and compare them with standard LLMs at matched scale.

\paragraph{Working memory in the brain.}
Language comprehension engages a left-lateralized language network, distinct from the domain-general multiple-demand (MD) network that supports working memory and cognitive control across tasks \citep{duncan2010multiple,fedorenko2013broad}. The MD network responds weakly during passive comprehension \citep{diachek2020domain}, and syntactic memory cost in naturalistic reading is tracked by the language network rather than by MD regions \citep{shain2022robust}. If memory-constrained models capture human working-memory limits, the MD network is a natural place to look for an advantage.

\paragraph{Representational similarity.}
Encoding models measure what a learned mapping can recover from a representation, and their capacity grows with representation width. RSA compares the geometry of stimuli in model and human spaces directly \citep{kriegeskorte2008rsa}. We use a held-out regression of the human dissimilarity matrix on the model one beyond lexical confounds, with stimulus-wise cross-validation \citep{schutt2023statistical} and standardized features \citep{timkey2021allbark}.
